\section{New Heavens and a New Earth: The Renewal of Creation}

God did not redeem the soul while abandoning the body, nor does He redeem humanity while casting the material universe onto a scrap heap. In Catholic eschatology, history ends not in the annihilation of creation, but in its cleansing and transfiguration into a home fit for the risen sons of God.

Through Isaias the Lord promises: ``For behold I create new heavens, and a new earth: and the former things shall not be in remembrance, and they shall not come upon the heart'' (Isa 65:17), sealing its permanence at the close of the book: ``For as the new heavens, and the new earth, which I will make to stand before me, saith the Lord: so shall your seed stand, and your name'' (Isa 66:22). In Romans 8, St.\ Paul sweeps the entire cosmos into the birth-pangs of redemption:

\begin{quote}
For I reckon that the sufferings of this time are not worthy to be compared with the glory to come that shall be revealed in us. For the expectation of the creature waiteth for the revelation of the sons of God. For the creature was made subject to vanity: not willingly, but by reason of him that made it subject, in hope. Because the creature also itself shall be delivered from the servitude of corruption, into the liberty of the glory of the children of God. For we know that every creature groaneth and travaileth in pain, even till now. (Rom 8:18--22)
\end{quote}

St.\ Peter repeats the promise after foretelling the purifying fire of judgment: ``But we look for new heavens and a new earth according to his promises, in which justice dwelleth'' (2 Pet 3:13). At the summit of Sacred Scripture, St.\ John beholds its fulfillment:

\begin{quote}
And I saw a new heaven and a new earth. For the first heaven and the first earth was gone: and the sea is now no more. And I, John, saw the holy city, the new Jerusalem, coming down out of heaven from God, prepared as a bride adorned for her husband. And I heard a great voice from the throne, saying: Behold the tabernacle of God with men: and he will dwell with them. And they shall be his people: and God himself with them shall be their God. And God shall wipe away all tears from their eyes: and death shall be no more. Nor mourning, nor crying, nor sorrow shall be any more, for the former things are passed away. And he that sat on the throne, said: Behold, I make all things new. (Apoc 21:1--5)
\end{quote}

That city needs neither stone temple nor sun or moon, ``for the Lord God Almighty is the temple thereof, and the Lamb,'' and ``the glory of God hath enlightened it: and the Lamb is the lamp thereof'' (Apoc 21:22--23). From the throne flows ``a river of water of life, clear as crystal,'' bordered by the tree of life bearing twelve fruits and leaves ``for the healing of the nations'' (Apoc 22:1--2). Every curse of Eden is lifted: ``And they shall see his face: and his name shall be on their foreheads. And night shall be no more \dots\ because the Lord God shall enlighten them. And they shall reign for ever and ever'' (Apoc 22:3--5).

Confronting second-century Gnostics who held the visible cosmos to be a flawed accident doomed to vanish, St.\ Irenaeus of Lyons laid down in \textit{Against Heresies} V.36.1 the principle that still governs Catholic doctrine:

\begin{quote}
For since there are real men, so must there also be a real establishment, that they vanish not away among non-existent things, but progress among those which have an actual existence. For neither is the substance nor the essence of the creation annihilated (for faithful and true is He who has established it), but ``the fashion of the world passeth away;'' that is, those things among which transgression has occurred, since man has grown old in them. \dots\ But when this fashion of things passes away, and man has been renewed, and flourishes in an incorruptible state, so as to preclude the possibility of becoming old, then there shall be the new heaven and the new earth, in which the new man shall remain continually, always holding fresh converse with God.\endnote{St.\ Irenaeus of Lyons, \textit{Against Heresies} V.36.1, trans.\ Roberts and Rambaut, ANF vol.\ 1.}
\end{quote}

Drawing on 1 Cor 7:31, Irenaeus distinguishes the enduring \textit{substance} of creation from its passing \textit{fashion} scarred by sin. At the close of the chapter (V.36.2--3), the redeemed ``ascend through the Spirit to the Son, and through the Son to the Father,'' until the Son hands the kingdom to the Father ``that God may be all in all'' (1 Cor 15:28).

In Book XX of \textit{The City of God} (XX.16--17), St.\ Augustine pairs physical renewal with sober restraint. Passing through the universal fire of the Last Day, ``the qualities of the corruptible elements which suited our corruptible bodies shall utterly perish, and our substance shall receive such qualities as shall, by a wonderful transmutation, harmonize with our immortal bodies, so that, as the world itself is renewed to some better thing, it is fitly accommodated to men, themselves renewed in their flesh to some better thing.''\endnote{St.\ Augustine, \textit{The City of God} XX.16, trans.\ Dods.} On ``there shall be no more sea,'' Augustine leaves open whether the physical sea is transformed or whether ``sea'' signifies the turbulent restlessness of mortal history. In Chapter 17 he rejects every attempt to flatten the New Jerusalem into an earthly utopia: in this age, ``the holier a man is, and the fuller of holy desire, so much the more abundant is the tearfulness of his supplication.'' Tears cease only when death itself is swallowed up.

Three centuries later, the Venerable Bede wove Augustine's vision into his \textit{Explanation of the Apocalypse} on Chapters 21 and 22.\endnote{Venerable Bede, \textit{Explanation of the Apocalypse}, Book III, on Chapters 21 and 22, trans.\ Edward Marshall (Oxford and London: James Parker and Co., 1878), pp.\ 142--169.} On Apoc 21:1--4, Bede writes that ``a heavenly incorruption will refine their bodies, and the sight of the eternal King will feed their mind''; the Bride now betrothed in faith is brought to an immortal marriage (21:9); the precious stone is Christ (21:11); the twelve gates are the Apostles (21:12); the city stands foursquare to signify ``the solidity of invincible truth'' (21:16). On the crystal river of Chapter 22, Bede joins sacrament to glory: ``In the river of life which flows in the midst of the city, there is no longer indicated the celebration of Baptism, but the fruit of that sacrament is now revealed. For in the present time the Church sows in the Spirit, that it may then reap eternal life.''\endnote{Venerable Bede, \textit{Explanation of the Apocalypse}, Book III, on Ch.\ 22, vv.\ 1--5, trans.\ Marshall, pp.\ 165--167.} Beside it stands the tree of life---Christ and the glory of the Cross---whose leaves signify ``a perpetual song'' for the healing of the nations; and ``night shall be no more'' because ``there no exhortation of the prophets, no preaching of the law, are necessary, which are now said to be `a lamp shining in a dark place'; seeing that in the contemplation of God all things are fulfilled.''

In the thirteenth century, the Thomistic \textit{Supplement} (q.\ 91, aa.\ 1, 3) explains why irrational creation shares in this glory: because corporeal things were made for man, ``the dwelling should befit the dweller,'' and ``that man's desire be satisfied the universe must needs also be made better'' (\textit{Supplement}, q.\ 91, a.\ 1, sed contra).\endnote{\textit{Supplement}, q.\ 91, aa.\ 1, 3 (English Dominicans, 1922).} Freed from the task of feeding mortal bodies, the cosmos serves a higher purpose:

\begin{quote}
Those bodies also will need to receive a greater inflow from the Divine goodness than now, not indeed so as to change their species, but so as to add a certain perfection of glory: and such will be the renewal of the world. Wherefore at the one same time, the world will be renewed, and man will be glorified. \dots\ The renewal of the world is directed to the end that, after this renewal has taken place, God may become visible to man by signs so manifest as to be perceived as it were by his senses.\endnote{\textit{Supplement}, q.\ 91, a.\ 1 and a.\ 3.}
\end{quote}

Even while teaching that the heavenly bodies will be heightened in brightness, the \textit{Supplement} keeps due humility (q.\ 91, a.\ 3): ``to what degree and in what way this betterment will take place is known to Him alone Who will bring it about.''

The \textit{Catechism of the Catholic Church} carries that exact poise forward: at the consummation of history, after the universal judgment, the just will reign forever with Christ, glorified in body and soul, and the material universe itself will be renewed (CCC 1042).\endnote{\textit{Catechism of the Catholic Church}, \S\S 1042, 1048 (citing Vatican II, \textit{Lumen gentium} 48 and \textit{Gaudium et spes} 39).} Echoing the Second Vatican Council's \textit{Gaudium et spes} (\S 39), it reminds the faithful that neither the moment of the earth's consummation nor the exact manner of cosmic transformation has been revealed; yet even as the sin-marred figure of this present world passes away, God is preparing a new dwelling and a new earth marked by justice, where happiness will fill and surpass every longing for peace arising in the human heart (CCC 1048).
